% Generated by tools/build_new_dists_anthology_sections.py; do not edit by hand.

\section{Symmetric, spin-cover, Lie, Foulkes, and rooted-tree groups}\label{proof:new-dist:symmetric_spin_lie_tree}

\resultbox{\textbf{Canonical testing artifacts.}\par\smallskip Exact backend: \path{lambda_distributions/dists/symmetric_spin_lie_tree_sigma_mgfs/verification.py}.\par Regression: \path{lambda_distributions/dists/symmetric_spin_lie_tree_sigma_mgfs/test_verification.py}.\par Marimo: \path{marimo_notebooks/symmetric_spin_lie_tree.py}.}

\subsection*{Scope and outcome}

We prove the generalized finite-group Molien identity and specialize it to
ordinary symmetric-group modules, genuine modules of a Schur double cover,
Foulkes permutation modules, multilinear free-Lie modules, iterated local
wreath products on regular rooted trees, and automorphism groups of arbitrary
finite rooted trees.  The rooted-tree representation is always the natural
permutation representation on level-$h$ leaves.  A reproducible marimo
notebook constructs every matrix group used in the audit and compares direct
fixed spaces or configuration orbits with the proposed character, Frobenius,
Witt, and cycle-index formulas.  The test suite performs $2{,}036$ comparisons;
all pass.  The finite formulas are separated from their asymptotic
consequences: fixed-tail Specht modules have Poisson-chaos limits, whereas raw
level-transitive leaf series cannot converge coefficientwise because their
pair coefficient is at least $h+1$.

\subsection{The common target and generalized Molien theorem}

Let $G$ be finite, let $\rho:G\to GL(V)$ be a complex representation with
character $\chi_V$, and let $\tau=(\tau_1,\ldots,\tau_s)$.  Put
\[
 W_\tau(V)=\bigotimes_{j=1}^s\Sym^{\tau_j}V,
 \qquad
 a_\tau(G,V)=\dim W_\tau(V)^G.
\]
The $\sigma$-MGF is the symmetric series
\[
 \MGF_{G,V}=\sum_\tau a_\tau(G,V)m_\tau.
\]

\begin{theorem}[generalized Molien identity]\label{nd:symmetric_spin_lie_tree:thm:molien}
For every finite group and finite-dimensional complex representation,
\begin{equation}\label{nd:symmetric_spin_lie_tree:eq:molien}
 \boxed{
 \MGF_{G,V}
 =\frac1{|G|}\sum_{g\in G}
 \exp\!\left(\sum_{r\ge1}\frac{p_r}{r}\chi_V(g^r)\right)
 =\sum_\tau\dim W_\tau(V)^Gm_\tau.}
\end{equation}
Equivalently, in an alphabet $\mathbf t=(t_1,t_2,\ldots)$,
\begin{equation}\label{nd:symmetric_spin_lie_tree:eq:det-molien}
 \MGF_{G,V}(\mathbf t)
 =\frac1{|G|}\sum_{g\in G}\prod_i
 \det(I-t_i\rho(g))^{-1}.
\end{equation}
\end{theorem}

\begin{proof}
On $W_\tau(V)$, the Reynolds operator
\[
 R_\tau=\frac1{|G|}\sum_{g\in G}\rho_{W_\tau}(g)
\]
is an idempotent with image $W_\tau(V)^G$.  Hence its trace is
$a_\tau(G,V)$.  For any matrix $A$,
\[
 \sum_{q\ge0}\Tr(\Sym^qA)t^q=\det(I-tA)^{-1}.
\]
Multiplying this identity over the tensor factors and averaging proves
\eqref{nd:symmetric_spin_lie_tree:eq:det-molien}.  Finally,
\[
 -\log\det(I-tA)=\sum_{r\ge1}\frac{t^r}{r}\Tr(A^r)
\]
and $\Tr(\rho(g)^r)=\chi_V(g^r)$ give \eqref{nd:symmetric_spin_lie_tree:eq:molien}.
\end{proof}

Thus a character table without power maps is not sufficient: the exponent
uses $\chi_V(g^r)$ for every $r$.  For a permutation $g$ of cycle type
$\mu=1^{m_1}2^{m_2}\cdots$, the $r$th power has type $\mu^{[r]}$ satisfying
\begin{equation}\label{nd:symmetric_spin_lie_tree:eq:power-map}
 \boxed{m_j(\mu^{[r]})=
 \sum_{\substack{d\mid r\\(j,r/d)=1}}d\,m_{jd}(\mu).}
\end{equation}
Indeed, a cycle of length $jd$ splits under the $r$th power into $d$ cycles
of length $j$ precisely when $d=(jd,r)$, which is equivalent to the displayed
divisibility and coprimality conditions.

\subsection{Ordinary symmetric-group modules}

\subsubsection{Frobenius master formula}

For an $S_n$-module $V_n$ with Frobenius characteristic $f_n=\ch(V_n)$,
the standard identity $\chi_{V_n}(\mu)=z_\mu[p_\mu]f_n$ and
Theorem~\ref{nd:symmetric_spin_lie_tree:thm:molien} give
\begin{equation}\label{nd:symmetric_spin_lie_tree:eq:frob-master}
 \boxed{
 \MGF_{V_n}
 =\sum_{\mu\vdash n}\frac1{z_\mu}
 \exp\!\left(
 \sum_{r\ge1}\frac{p_r}{r}
 z_{\mu^{[r]}}[p_{\mu^{[r]}}]f_n
 \right).}
\end{equation}
For a Specht module $S^\lambda$, this specializes to
\begin{equation}\label{nd:symmetric_spin_lie_tree:eq:specht}
 \boxed{
 \MGF_{S_n,S^\lambda}
 =\sum_{\mu\vdash n}\frac1{z_\mu}
 \exp\!\left(\sum_{r\ge1}\frac{p_r}{r}
 \chi^\lambda(\mu^{[r]})\right).}
\end{equation}
The formula is exact for every $n$ and $\lambda$; character evaluation can be
performed by Murnaghan--Nakayama or any equivalent character algorithm.

Every irreducible $S_n$-module is self-dual and admits an invariant symmetric
bilinear form.  Therefore, for a nontrivial irreducible,
\[
 [m_{(1)}]=0,
 \qquad [m_{(2)}]=[m_{(1,1)}]=1.
\]
The fourth multilinear coefficient already sees Kronecker complexity:
\[
 [m_{(1,1,1)}]=g(\lambda,\lambda,\lambda),
 \qquad
 [m_{(1,1,1,1)}]=\sum_{\nu\vdash n}g(\lambda,\lambda,\nu)^2.
\]
Consequently \eqref{nd:symmetric_spin_lie_tree:eq:specht} is a complete exact formula, but not a claim
that arbitrary Kronecker multiplicities have a simpler closed form.

\subsubsection{Hooks, two rows, and the matrix audit}

Let $C_j(\sigma)$ count $j$-cycles.  Since
$S^{(n-k,1^k)}\cong\bigwedge^kS^{(n-1,1)}$, its character is
\begin{equation}\label{nd:symmetric_spin_lie_tree:eq:hook}
 \boxed{Q_{1^k}(C)=[u^k]\frac1{1+u}
 \prod_{j\ge1}(1-(-u)^j)^{C_j}.}
\end{equation}
This follows by deleting the trivial eigenline from the natural permutation
module and taking the generating function of exterior-power characters.

For $n\ge2k$, let
\[
 A_k(C)=[u^k]\prod_{j\ge1}(1+u^j)^{C_j}.
\]
This is the number of fixed $k$-subsets.  Young's rule gives
\[
 \C\!\left[\binom{[n]}k\right]
 \cong\bigoplus_{j=0}^kS^{(n-j,j)},
\]
and subtraction of consecutive characters proves
\begin{equation}\label{nd:symmetric_spin_lie_tree:eq:two-row}
 \boxed{\chi^{(n-k,k)}=A_k-A_{k-1}.}
\end{equation}

The verifier constructs the standard module on the integral basis
$e_i-e_n$, takes exterior minors for the hook cases, and realizes
$S^{(n-2,2)}$ as the vertex-sum-zero subspace of the two-subset permutation
module.  It checks the traces of powers $1\le r\le4$ for every matrix against
\eqref{nd:symmetric_spin_lie_tree:eq:hook} or \eqref{nd:symmetric_spin_lie_tree:eq:two-row}.  It then constructs
$W_\tau(V)$ for $\tau=(1),(2),(1,1),(3)$, finds the common generator-fixed
space, and compares its dimension with the full Molien average.

\begin{center}
\begin{tabular}{@{}lrrrrr@{}}\toprule
module&dim.&$[m_1]$&$[m_2]$&$[m_{11}]$&$[m_3]$\\\midrule
$S^{(3,1)}$ of $S_4$&3&0&1&1&1\\
$S^{(3,1,1)}$ of $S_5$&6&0&1&1&0\\
$S^{(2,2)}$ of $S_4$&2&0&1&1&1\\
$S^{(3,2)}$ of $S_5$&5&0&1&1&1\\\bottomrule
\end{tabular}
\end{center}
All $1{,}152$ character-power comparisons and all $16$ fixed-space/Molien
comparisons pass.  Numerical rank residuals are at most $1.3\times10^{-15}$;
the matrices and character calculations themselves have integral entries.

\subsubsection{Fixed-tail asymptotics and the boundary of the claim}

Fix a partition $\lambda$ and put
$\lambda[n]=(n-|\lambda|,\lambda_1,\lambda_2,\ldots)$.  In the stable range,
$\chi^{\lambda[n]}$ is a fixed character polynomial
$Q_\lambda(C_1,C_2,\ldots)$.  Let independent
$Z_j\sim\operatorname{Poisson}(1/j)$ and define
\begin{equation}\label{nd:symmetric_spin_lie_tree:eq:poisson-power}
 Z_j^{[r]}=\sum_{\substack{d\mid r\\(j,r/d)=1}}dZ_{jd}.
\end{equation}
Then, coefficientwise in bounded total degree,
\begin{equation}\label{nd:symmetric_spin_lie_tree:eq:fixed-tail-limit}
 \boxed{
 \MGF_{\lambda,\infty}
 =\E\exp\!\left(
 \sum_{r\ge1}\frac{p_r}{r}
 Q_\lambda(Z_1^{[r]},Z_2^{[r]},\ldots)
 \right).}
\end{equation}

\begin{proof}
In total degree at most $D$, the exponential uses only $r\le D$ and finitely
many cycle variables in each character polynomial.  Formula
\eqref{nd:symmetric_spin_lie_tree:eq:power-map} expresses the cycle counts of $\sigma^r$ in those of
$\sigma$.  The finite collection of cycle counts of a uniform permutation
converges jointly, with all fixed moments, to the independent $Z_j$.
Substitution and moment convergence prove \eqref{nd:symmetric_spin_lie_tree:eq:fixed-tail-limit}.
\end{proof}

This is a polynomial Poisson-chaos law, not generally an ordinary Poisson
law.  It applies to fixed hooks and fixed two-row tails through
\eqref{nd:symmetric_spin_lie_tree:eq:hook}--\eqref{nd:symmetric_spin_lie_tree:eq:two-row}.  It does not apply to rectangles or
diagrams with two growing dimensions.  For those shapes \eqref{nd:symmetric_spin_lie_tree:eq:specht}
remains exact, while normalized character asymptotics alone do not imply a
raw $\sigma$-MGF limit; boundedness of the fourth coefficient is already a
necessary test.

\subsection{Schur double covers and genuine modules}

Let
\[
 1\longrightarrow\{1,z\}\longrightarrow\widetilde S_n
 \longrightarrow S_n\longrightarrow1
\]
be a Schur double cover, and let $D$ be a genuine module, so $z$ acts by
$-I$.  Theorem~\ref{nd:symmetric_spin_lie_tree:thm:molien}, grouped by conjugacy classes, gives
\begin{equation}\label{nd:symmetric_spin_lie_tree:eq:schur-cover}
 \boxed{
 \MGF_{\widetilde S_n,D}
 =\sum_{[\widetilde g]}
 \frac1{|C_{\widetilde S_n}(\widetilde g)|}
 \exp\!\left(\sum_{r\ge1}\frac{p_r}{r}
 \chi_D(\widetilde g^r)\right).}
\end{equation}
Pairing the summands $g$ and $zg$ shows immediately that
\begin{equation}\label{nd:symmetric_spin_lie_tree:eq:central-parity}
 \boxed{[m_\tau]\MGF_{\widetilde S_n,D}=0
 \quad\text{when }|\tau|\text{ is odd}.}
\end{equation}

If $D\otimes D\cong\bigoplus_\eta m_\eta U_\eta$ and the required
self-duality identifies fourth invariants with
$\operatorname{End}_{\widetilde S_n}(D\otimes D)$, Schur's lemma gives
\begin{equation}\label{nd:symmetric_spin_lie_tree:eq:spin-fourth}
 \dim(D^{\otimes4})^{\widetilde S_n}=\sum_\eta m_\eta^2.
\end{equation}
Without that self-duality, the always-valid statement pairs multiplicities
of dual constituents instead.  This qualification is essential when the
chosen genuine constituent is not self-dual.

For the audit, Hermitian Clifford matrices $\gamma_i$ satisfy
$\gamma_i\gamma_j+\gamma_j\gamma_i=2\delta_{ij}I$.  The lifts
\[
 u_i=\frac{\gamma_i-\gamma_{i+1}}{\sqrt2},\qquad 1\le i<4,
\]
generate the Pin preimage of $S_4$, an order-$48$ matrix group containing
$-I$.  The represented module has dimension four.  Direct generator-fixed
spaces and \eqref{nd:symmetric_spin_lie_tree:eq:schur-cover} give
\begin{center}
\begin{tabular}{@{}crrrrr@{}}\toprule
$\tau$&$(1)$&$(2)$&$(1,1)$&$(3)$&$(1,1,1,1)$\\\midrule
coefficient&0&1&2&0&16\\\bottomrule
\end{tabular}
\end{center}
All five comparisons pass, the odd rows vanish as predicted, and the maximum
Clifford-relation residual is zero at machine precision.  This test verifies
the universal formula and central selection rule for a genuine matrix module;
it does not purport to tabulate every strict-partition spin character.

\subsection{Foulkes permutation modules}

Let
\[
 H^{(a^b)}=\Ind_{S_a\wr S_b}^{S_{ab}}\mathbf1.
\]
It is the permutation module on partitions of $[ab]$ into $b$ unlabeled
blocks of size $a$, and its Frobenius characteristic is
\[
 \ch H^{(a^b)}=h_b[h_a].
\]
Substitution in \eqref{nd:symmetric_spin_lie_tree:eq:frob-master} proves
\begin{equation}\label{nd:symmetric_spin_lie_tree:eq:foulkes}
 \boxed{
 \MGF_{H^{(a^b)}}=
 \sum_{\mu\vdash ab}\frac1{z_\mu}
 \exp\!\left(
 \sum_{r\ge1}\frac{p_r}{r}
 z_{\mu^{[r]}}[p_{\mu^{[r]}}]h_b[h_a]
 \right).}
\end{equation}

The pair coefficient has a concrete orbit interpretation.  Given two block
partitions, label their blocks temporarily and set
$A_{ij}=|B_i\cap C_j|$.  Every row and column sum is $a$.  Relabeling the two
block systems permutes rows and columns, and two pairs lie in the same
$S_{ab}$-orbit exactly when the tables are so related.  Therefore
\begin{equation}\label{nd:symmetric_spin_lie_tree:eq:contingency}
 \boxed{[m_{(1,1)}]\MGF_{H^{(a^b)}}
 =\#\{A\in\mathbb Z_{\ge0}^{b\times b}:
 \text{all margins are }a\}/(S_b\times S_b).}
\end{equation}

The verification lists the block partitions, constructs every zero-one
permutation matrix, and independently expands
\[
 h_b[h_a]=\sum_{\nu\vdash b}\frac1{z_\nu}
 \prod_{r\in\nu}\left(\sum_{\lambda\vdash a}
 \frac{p_{r\lambda}}{z_\lambda}\right).
\]
Thus every matrix trace is compared with a power-sum coefficient, not with a
second traversal of the same action.  The results are
\begin{center}
\begin{tabular}{@{}lrrrrrr@{}}\toprule
module&dim.&$[m_1]$&$[m_2]$&$[m_{11}]$&$[m_3]$&$[m_{21}]$\\\midrule
$H^{(2^2)}$ of $S_4$&3&1&2&2&3&4\\
$H^{(2^3)}$ of $S_6$&15&1&3&3&8&12\\\bottomrule
\end{tabular}
\end{center}
All $744$ character comparisons and ten direct-orbit/Molien comparisons pass;
the contingency-table counts are respectively $2$ and $3$.  Their growth in
the standard Foulkes regimes obstructs an unnormalized coefficientwise
limit.

\subsection{The multilinear free-Lie representation}

The degree-$n$ multilinear free-Lie module has Frobenius characteristic
\begin{equation}\label{nd:symmetric_spin_lie_tree:eq:witt}
 \boxed{L_n=\ch\Lie_n=\frac1n\sum_{d\mid n}\mu(d)p_d^{n/d}.}
\end{equation}
Formula \eqref{nd:symmetric_spin_lie_tree:eq:frob-master} therefore yields
\begin{equation}\label{nd:symmetric_spin_lie_tree:eq:lie-smg}
 \boxed{
 \MGF_{\Lie_n}=
 \sum_{\mu\vdash n}\frac1{z_\mu}
 \exp\!\left(
 \sum_{r\ge1}\frac{p_r}{r}
 z_{\mu^{[r]}}[p_{\mu^{[r]}}]L_n
 \right).}
\end{equation}
Power sums are orthogonal with
$\langle p_\lambda,p_\lambda\rangle=z_\lambda$.  Distinct divisors in
\eqref{nd:symmetric_spin_lie_tree:eq:witt} give distinct cycle types, hence
\begin{equation}\label{nd:symmetric_spin_lie_tree:eq:lie-norm}
 \boxed{[m_{(1,1)}]\MGF_{\Lie_n}
 =\langle L_n,L_n\rangle
 =\frac1{n^2}\sum_{d\mid n}\mu(d)^2d^{n/d}(n/d)!.}
\end{equation}
The $d=1$ term gives the lower bound $n!/n^2$, so the raw series cannot have
a coefficientwise limit as $n\to\infty$.  A sign twist, such as that appearing
in top partition-lattice homology, leaves this norm unchanged.

For an independent matrix construction, fix $x_n$ and use the bracket basis
\[
 [\cdots[[x_n,x_{\sigma(1)}],x_{\sigma(2)}],\ldots,
 x_{\sigma(n-1)}],\qquad \sigma\in S_{n-1}.
\]
Expanding brackets in the multilinear associative word basis gives integral
matrices for relabeling by $S_n$.  For $n=3,4$, all $30$ matrix characters
match \eqref{nd:symmetric_spin_lie_tree:eq:witt}.  Direct fixed spaces and Molien averages give
\begin{center}
\begin{tabular}{@{}lrrrrr@{}}\toprule
module&dim.&$[m_1]$&$[m_2]$&$[m_{11}]$&$[m_3]$\\\midrule
$\Lie_3$&2&0&1&1&1\\
$\Lie_4$&6&0&2&2&2\\\bottomrule
\end{tabular}
\end{center}
The pair values $1,2$ agree exactly with \eqref{nd:symmetric_spin_lie_tree:eq:lie-norm}.

\subsection{Rooted-tree groups in their natural leaf representations}

\subsubsection{Elementwise wreath recursion}

Let $T_{d,h}$ be the rooted $d$-ary tree of height $h$, and let $X_h$ be its
level-$h$ leaves.  Every automorphism has wreath form
\[
 g=(g_1,\ldots,g_d)\pi,\qquad \pi\in S_d.
\]
For a cycle $C=(i_1\cdots i_\ell)$ of $\pi$, define the section product
$g_C=g_{i_\ell}\cdots g_{i_1}$.  Put
\[
 D_h(g;\mathbf t)=\prod_i
 \det(I-t_iP_g\mid\C[X_h])^{-1}.
\]

\begin{proposition}[elementwise recursion]\label{nd:symmetric_spin_lie_tree:prop:elementwise-tree}
\begin{equation}\label{nd:symmetric_spin_lie_tree:eq:elementwise-tree}
 \boxed{D_h(g;\mathbf t)=
 \prod_{C\in\Cyc(\pi)}
 D_{h-1}(g_C;t_1^{|C|},t_2^{|C|},\ldots).}
\end{equation}
\end{proposition}

\begin{proof}
If $g_C$ has an $m$-cycle on the lower leaves, one circuit around a top cycle
of length $|C|$ produces a single $m|C|$-cycle.  Its determinant factor is
$1-t_i^{m|C|}$.  Multiplication over lower cycles, top cycles, and alphabets
proves \eqref{nd:symmetric_spin_lie_tree:eq:elementwise-tree}.
\end{proof}

For any finite level quotient $G_h$,
\begin{equation}\label{nd:symmetric_spin_lie_tree:eq:tree-average}
 \boxed{\MGF_{G_h,X_h}=\frac1{|G_h|}\sum_{g\in G_h}D_h(g;\mathbf t).}
\end{equation}
This is exact for every supplied finite automaton or branch quotient once its
allowed wreath recursions are specified.  Level transitivity alone does not
determine the full finite series; the actual quotient and its sections remain
part of the data.

\subsubsection{Iterated local wreath products}

Let $L\le S_d$ be transitive and define $G_0=1$ and
$G_{h+1}=G_h\wr L$ in the rooted-tree action.  Write
$\Phi_h=\MGF_{G_h,X_h}$.  Independence and uniformity of section products in
the disjoint cycles of $\pi\in L$ give
\begin{align}
 \Phi_0(\mathbf t)&=\prod_i(1-t_i)^{-1},\label{nd:symmetric_spin_lie_tree:eq:phi-zero}\\
 \Phi_{h+1}(\mathbf t)
 &=\frac1{|L|}\sum_{\pi\in L}
 \prod_{C\in\Cyc(\pi)}
 \Phi_h(t_1^{|C|},t_2^{|C|},\ldots).\label{nd:symmetric_spin_lie_tree:eq:phi-recurrence}
\end{align}
For the regular cyclic action $L=C_p$, this becomes
\begin{equation}\label{nd:symmetric_spin_lie_tree:eq:cyclic-tree}
 \boxed{\Phi_{h+1}=\frac1p\left(
 \Phi_h(\mathbf t)^p+(p-1)\Phi_h(\mathbf t^p)\right).}
\end{equation}
This is the natural leaf action of a Sylow $p$-subgroup of $S_{p^h}$.
For $L=S_d$, the cycle index gives
\begin{equation}\label{nd:symmetric_spin_lie_tree:eq:full-tree}
 \boxed{\Phi_{h+1}=
 \sum_{\lambda\vdash d}\frac1{z_\lambda}
 \prod_{\ell\ge1}\Phi_h(\mathbf t^\ell)^{m_\ell(\lambda)}.}
\end{equation}

\subsubsection{Fixed leaves, branching moments, and the no-limit theorem}

For uniform $g\in G_h$, let $X_h=|\Fix_{X_h}(g)|$.  If
$Y=|\Fix_{[d]}(\pi)|$ for uniform $\pi\in L$, then
\begin{equation}\label{nd:symmetric_spin_lie_tree:eq:branching}
 X_{h+1}\overset d=\sum_{j=1}^{Y}X_h^{(j)},
\end{equation}
where the summands are independent copies of $X_h$.  Thus $X_h$ is the
generation-size process of a Galton--Watson process.  Transitivity and
Burnside's lemma imply
\[
 \E Y=1,
 \qquad \E Y^2=r_L:=\#(L\backslash[d]^2).
\]
Conditioning in \eqref{nd:symmetric_spin_lie_tree:eq:branching} proves by induction
\begin{equation}\label{nd:symmetric_spin_lie_tree:eq:second-moment}
 \boxed{\E X_h=1,
 \qquad \E X_h^2=1+h(r_L-1).}
\end{equation}
Since $[m_{(1,1)}]\Phi_h=\E X_h^2$, every nontrivial transitive local action
fails to have a raw coefficientwise limit.  In particular,
\begin{align*}
 L=C_p:&\quad [m_{(1,1)}]\Phi_h=1+(p-1)h,\\
 L=S_d:&\quad [m_{(1,1)}]\Phi_h=h+1.
\end{align*}
When $0<\operatorname{Var}Y=r_L-1<\infty$, the standard critical branching
estimate further gives
\[
 \Pr(X_h>0)\sim\frac{2}{(r_L-1)h}.
\]
Thus fixed leaves disappear in probability while the second moment diverges.

There is also a recurrence-free obstruction.

\begin{theorem}[universal level-transitive obstruction]\label{nd:symmetric_spin_lie_tree:thm:tree-obstruction}
If $G_h\le\Aut(T_{d,h})$ is transitive on $X_h$, then
\begin{equation}\label{nd:symmetric_spin_lie_tree:eq:universal-bound}
 \boxed{[m_{(1,1)}]\MGF_{G_h,X_h}
 =\#(G_h\backslash X_h^2)\ge h+1.}
\end{equation}
Consequently no level-transitive family of natural leaf representations has
a raw coefficientwise $\sigma$-MGF limit as $h\to\infty$.
\end{theorem}

\begin{proof}
For an ordered pair of leaves, every rooted-tree automorphism preserves the
depth of the lowest common ancestor.  Equality and the $h$ possible
divergence depths give $h+1$ nonempty invariant subsets of $X_h^2$, hence at
least that many orbits.  The orbit interpretation of the pair coefficient is
Burnside's lemma.
\end{proof}

This applies, under the explicit level-transitivity hypothesis, to the usual
finite quotients of many Grigorchuk, Gupta--Sidki, GGS, Basilica, and iterated
monodromy actions.  It asserts failure of raw moment convergence, not failure
of convergence in probability or of a suitably normalized theory.

\subsubsection{Exact iterated-group verification}

Every wreath element is materialized as a leaf permutation.  The direct route
extracts symmetric-power traces from its cycle structure.  The independent
route performs truncated multivariate polynomial arithmetic in
\eqref{nd:symmetric_spin_lie_tree:eq:phi-recurrence}.  Direct fixed-point squares and direct pair-orbit
search separately verify \eqref{nd:symmetric_spin_lie_tree:eq:second-moment} and
\eqref{nd:symmetric_spin_lie_tree:eq:universal-bound}.

\begin{center}
\begin{tabular}{@{}lrrrrrr@{}}\toprule
local action&$h$&leaves&$|G_h|$&direct $\E X_h^2$&formula&pair orbits\\\midrule
$C_2$&1&2&2&2&2&2\\
$C_2$&2&4&8&3&3&3\\
$C_2$&3&8&128&4&4&4\\
$C_3$&1&3&3&3&3&3\\
$C_3$&2&9&81&5&5&5\\
$S_3$&1&3&6&2&2&2\\
$S_3$&2&9&1296&3&3&3\\\bottomrule
\end{tabular}
\end{center}
For each row, the six coefficients indexed by
$(1),(2),(1,1),(3),(2,1),(2,2)$ agree exactly with the recurrence.  All
$42$ coefficient comparisons and all moment/orbit checks pass.

\subsection{Arbitrary finite and random rooted trees}

Let $T$ be a finite rooted tree.  Group the children of its root by rooted
isomorphism type $U$, with multiplicity $m_U$.  Then
\[
 \Aut(T)\cong\prod_{[U]}\bigl(\Aut(U)\wr S_{m_U}\bigr).
\]
For the symmetric-group cycle index
\[
 Z_{S_m}(A_1,\ldots,A_m)=\frac1{m!}\sum_{\pi\in S_m}
 \prod_{C\in\Cyc(\pi)}A_{|C|},
\]
the direct-product and wreath-product determinant rules prove
\begin{equation}\label{nd:symmetric_spin_lie_tree:eq:arbitrary-tree}
 \boxed{
 \Phi_T(\mathbf t)=\prod_{[U]}Z_{S_{m_U}}
 \bigl(\Phi_U(\mathbf t),\Phi_U(\mathbf t^2),\ldots,
 \Phi_U(\mathbf t^{m_U})\bigr).}
\end{equation}
For a random tree $\mathcal T$, conditioning on the realized branch types
turns \eqref{nd:symmetric_spin_lie_tree:eq:arbitrary-tree} into the corresponding recursive
distributional identity.  This is a finite-height computational recursion.
It does not by itself imply a universal height-infinity law: asymptotics for
the scalar $|\Aut(\mathcal T)|$ do not control every orbit coefficient of the
much richer random series.

The verifier constructs three non-regular trees with four or five leaves and
automorphism-group orders $4,8,12$.  For each it compares direct coefficients
at $(1),(2),(1,1),(3),(2,1)$ with \eqref{nd:symmetric_spin_lie_tree:eq:arbitrary-tree}.  All $15$
comparisons pass exactly.

\subsection{Verification architecture and reproduction}

The two computational paths are intentionally different.
\begin{enumerate}
\item \textbf{Matrix/fixed-space path.}  Specht and free-Lie matrices are
integral.  Clifford matrices generate the Schur cover.  Symmetric-square,
symmetric-cube, and tensor matrices are formed explicitly, and the common
fixed space of group generators is computed.  The maximum reported numerical
nullspace residual is $1.7\times10^{-15}$.
\item \textbf{Character/formula path.}  Every group element is averaged using
power traces and Newton's identity.  Hook/two-row character polynomials,
power-sum plethysm for $h_b[h_a]$, the Witt characteristic, and truncated
wreath cycle indices supply the family-specific formulas.
\item \textbf{Literal orbit path.}  Foulkes and tree permutation modules are
also checked by enumerating multisets or ordered pairs and counting their
orbits, avoiding character algebra altogether.
\end{enumerate}

From the workspace root, run
\begin{verbatim}
.venv/bin/python -m \
  lambda_distributions.dists.symmetric_spin_lie_tree_sigma_mgfs.verification
.venv/bin/pytest -q \
  lambda_distributions/dists/symmetric_spin_lie_tree_sigma_mgfs/test_verification.py
.venv/bin/marimo edit \
  marimo_notebooks/symmetric_spin_lie_tree.py
\end{verbatim}
Join wrapped commands onto one line.  The noninteractive program reports
\begin{verbatim}
PASS: 2036 formula, character, matrix, and orbit checks
\end{verbatim}

\subsection{Conclusions and asymptotic classification}

\begin{center}
\small
\begin{tabularx}{\linewidth}{@{}>{\raggedright\arraybackslash}p{0.25\linewidth}
>{\raggedright\arraybackslash}p{0.27\linewidth}X@{}}\toprule
family&exact finite formula&asymptotic status\\\midrule
fixed-tail Specht&\eqref{nd:symmetric_spin_lie_tree:eq:specht}, \eqref{nd:symmetric_spin_lie_tree:eq:fixed-tail-limit}&explicit polynomial Poisson chaos\\
hooks and two rows&\eqref{nd:symmetric_spin_lie_tree:eq:hook}, \eqref{nd:symmetric_spin_lie_tree:eq:two-row}&explicit in the fixed-tail regime\\
growing diagrams&\eqref{nd:symmetric_spin_lie_tree:eq:specht}&raw limit not supplied by normalized character theory\\
genuine $2.S_n$ modules&\eqref{nd:symmetric_spin_lie_tree:eq:schur-cover}, \eqref{nd:symmetric_spin_lie_tree:eq:central-parity}&normalized spin theory does not alone control raw moments\\
Foulkes&\eqref{nd:symmetric_spin_lie_tree:eq:foulkes}, \eqref{nd:symmetric_spin_lie_tree:eq:contingency}&pair-orbit growth obstructs standard raw limits\\
$\Lie_n$&\eqref{nd:symmetric_spin_lie_tree:eq:lie-smg}, \eqref{nd:symmetric_spin_lie_tree:eq:lie-norm}&$[m_{11}]\ge n!/n^2$, so no raw limit\\
iterated local tree groups&\eqref{nd:symmetric_spin_lie_tree:eq:phi-recurrence}&critical branching; no raw leaf limit\\
arbitrary finite trees&\eqref{nd:symmetric_spin_lie_tree:eq:arbitrary-tree}&finite recursion; random infinite-height law is model-dependent\\\bottomrule
\end{tabularx}
\end{center}

The central distinction is between an exact finite formula and a limiting
law.  All finite identities above follow from Reynolds averaging plus a
family-specific character or cycle recursion and are verified on explicit
matrix groups.  Only the fixed-tail Specht regime yields the stated positive
coefficientwise limit.  The Foulkes, free-Lie, and level-transitive tree
families have explicit low-degree coefficients that diverge, so their raw
series cannot converge even when individual normalized characters or fixed
points have useful limits.

\begin{thebibliography}{9}
\bibitem{nd:symmetric_spin_lie_tree:cef}
T.~Church, J.~Ellenberg, and B.~Farb,
\emph{FI-modules and stability for representations of symmetric groups},
\url{https://arxiv.org/abs/1204.4533}.

\bibitem{nd:symmetric_spin_lie_tree:spin}
S.~Matsumoto and P.~\'{S}niady,
\emph{Linear versus spin: representation theory of the symmetric groups},
\url{https://arxiv.org/abs/1811.10434}.

\bibitem{nd:symmetric_spin_lie_tree:foulkes}
C.~de Boeck,
\emph{On the decomposition of the Foulkes module},
\url{https://arxiv.org/abs/1207.6300}.

\bibitem{nd:symmetric_spin_lie_tree:lie}
S.~Sundaram,
\emph{Prime power variations of higher $\Lie_n$ modules},
\url{https://arxiv.org/abs/2107.06389}.

\bibitem{nd:symmetric_spin_lie_tree:img}
V.~Nekrashevych,
\emph{Iterated Monodromy Groups},
\url{https://arxiv.org/abs/math/0312306}.

\bibitem{nd:symmetric_spin_lie_tree:wreath}
A.~Skuratovskii,
\emph{Generalized iterated wreath products of cyclic groups and rooted trees correspondence},
\url{https://arxiv.org/abs/1409.0603}.

\bibitem{nd:symmetric_spin_lie_tree:randomtrees}
M.~Olsson and S.~Wagner,
\emph{The distribution of the number of automorphisms of random trees},
\url{https://arxiv.org/abs/2211.12801}.
\end{thebibliography}
